% ============================================================
% 第八章 结论与展望
% ============================================================

\chapter{结论与展望}

\section{研究结论}

本文围绕"基于充气结构的车载多形态飞行机器人系统（飞行雨伞）"这一课题，
完成了从概念设计、理论计算到样机制作与实验验证的全流程工程实现。
主要研究结论如下：

\textbf{（1）充气结构赋能形态主动重构的技术可行性已得到验证。}

以快递充气气柱为骨架材料，结合PLA 3D打印节点与TPU伞面，
本文设计并制作了一套能够在收纳态、飞行态、遮雨态之间主动切换的充气结构系统。
气柱骨架充气展开时间约% TODO-未填写\,s，气密性保持时间超过% TODO-未填写\,min，
排气收纳时间约% TODO-未填写\,s，三种形态切换的完整工作闭环得到实验验证。

\textbf{（2）充气伞形多旋翼系统的飞行可行性已得到验证。}

基于1230\,mm轴距正方形X型四旋翼布局，
配合SIRIUS 2306.5-1850KV电机与51366 MCK v3桨叶，
系统实现了稳定悬停飞行。
整机起飞质量1468\,g，悬停工作油门约26\%$\sim$30\%，
推重比达2.32（50\%油门），系留模式理论续航约40\,min，
各动力指标均满足设计需求。
转动惯量约为同类标准四旋翼的16倍（$I_{xx} \approx 0.196$\,kg$\cdot$m$^2$），
针对该特性的PID整定策略在实飞验证中% TODO-未填写（如：取得了良好效果，系统悬停姿态稳定）。

\textbf{（3）车载快速部署场景的工程方案得到验证。}

系统采用12\,V车载电源驱动充气泵，零改装集成于新能源汽车后备箱，
从取出到完成起飞的准备时间约% TODO-未填写\,s，
完整工作流程（收纳→展开→飞行→收纳）得到端到端验证，
"随车常备、快速部署"的使用范式具备工程可行性。

\textbf{（4）两种供电形态的互补设计满足不同场景需求。}

自载电池模式（6S 1500\,mAh）提供应急备用电源与短时功能验证能力；
系留供电模式（6S 10000\,mAh）通过地面有线供电实现长时悬停，
用户电源可随身携带于口袋或背包，实现双手解放。
两种模式插拔切换，相互补充，共同构成完整的能源管理方案。

\textbf{（5）伞面几何设计满足遮蔽功能需求。}

伞面水平投影面积0.471\,m$^2$（正方形，边长686\,mm），
伞面倾角约12.8°$\sim$13.4°，与传统雨伞设计经验高度吻合。
实测遮蔽面积约% TODO-未填写\,m$^2$，
满足设计指标（$\geq 0.40$\,m$^2$），防水遮蔽功能% TODO-未填写（如：通过验证）。


\section{创新点总结}

本文的主要创新点可归纳为三个层面：

\begin{enumerate}
    \item \textbf{充气结构作为飞行器形态主动重构核心手段}：
          首次将薄壁充气气柱引入多旋翼飞行器形态重构设计，
          以充气/排气驱动系统在紧凑收纳态与大面积飞行遮蔽态之间主动切换，
          提供了面积量级远大于刚性折叠方案的形态重构路径。

    \item \textbf{面向充气伞形载体的多旋翼布局专项设计}：
          针对方形充气伞体的大面积、气动阻力显著、转动惯量超大等特殊动力学特性，
          开展旋翼布局、推重比设计与PID整定的专项研究，
          验证了超大转动惯量（约标准四旋翼16倍）场景下的稳定悬停飞行可行性。

    \item \textbf{车载部署与个人出行服务的系统化方案}：
          将飞行机器人定位为"车载服务系统的延伸执行器"，
          提出涵盖收纳固定、快速充气展开、离车悬停服务与归位回收的完整车载部署方案，
          回应了新能源汽车时代"车载服务向用户物理空间延伸"的现实需求，
          在现有车载无人机研究中属于新的应用范式探索。
\end{enumerate}


\section{不足与局限}

受限于样机阶段的工程条件与时间约束，本系统在以下方面存在明显不足，
需在后续研究中加以改进：

\begin{itemize}
    \item \textbf{自主跟随能力缺失}：
          当前系统依赖操作者手动遥控实现跟随，用户体验与无人机自主性受限。
          真正的"无感知负担"体验需要计算机视觉目标检测与自主跟随控制系统的支撑，
          这是从样机走向产品的核心技术壁垒。

    \item \textbf{悬停精度有限}：
          在无GPS辅助的室内场景或GPS信号较弱的环境下，
          系统位置保持依赖操作者手动修正，悬停精度约$\pm 0.3\sim0.5$\,m，
          在风较大时对准用户头顶的难度增大。

    \item \textbf{充气结构结构强度有待提升}：
          当前样机使用快递气柱作为骨架材料，其截面形状不规则、
          节点连接刚度有限，在较强侧风（$>3$\,m/s）下骨架存在一定形变。
          后续应探索定制截面气柱或硬壳气柱方案，提升结构刚度。

    \item \textbf{噪声较大}：
          51366桨叶在飞行状态下噪声约% TODO-未填写\,dB，
          近距离（1.5\,m以内）使用时对用户构成一定噪声干扰，
          有待通过低噪声桨叶设计或降低旋翼转速加以改善。

    \item \textbf{安全性待进一步验证}：
          飞行器长期悬停于用户头顶，对飞控可靠性与失效模式的安全设计要求极高。
          当前样机在失控保护与机械安全结构方面仅完成了初步设计，
          距离实际产品的安全标准仍有较大距离，需要系统性的失效模式分析（FMEA）
          与冗余设计。
\end{itemize}


\section{未来工作展望}

基于本文的工程实践，后续研究工作可从以下四个方向展开：

\textbf{方向一：自主视觉跟随}

在飞行雨伞平台上集成轻量化计算模块（如NVIDIA Jetson Nano或手机芯片SoC），
结合单目/深度摄像头与目标检测算法（如YOLOv8-Nano），
实现对用户的实时目标识别、深度估计与跟踪控制，
使飞行雨伞能够在无人工干预的情况下自主跟随用户步行，
真正实现"双手解放、全程自动"的使用体验。

\textbf{方向二：低噪声旋翼系统优化}

研究低桨盘载荷（更大直径、更低转速）的旋翼方案，
配合主动降噪或旋翼声学优化设计，
将飞行雨伞的使用噪声控制在令人舒适的范围内（目标$\leq 65$\,dB @ 1\,m），
提升近人飞行的用户体验。

\textbf{方向三：结构强度与轻量化的协同优化}

探索定制截面充气气柱（如圆形截面高压气柱）与碳纤维骨架混合结构方案，
在保持折叠收纳特性的同时提升骨架刚度与抗侧风能力；
同时通过拓扑优化减轻PLA节点结构件质量，
进一步降低整机起飞质量，提升续航时间与动力裕量。

\textbf{方向四：车载一体化平台产品化}

设计专用车载底座（适配常见SUV后备箱导轨），
集成自动充气泵、快速展开机构与无线充电功能，
实现飞行雨伞从"工程样机"到"车载电子产品"的跨越，
最终探索将飞行雨伞作为新能源汽车生态配件的商业化路径。

\vspace{1em}
飞行雨伞作为一个将传统雨伞重新定义的创新概念，
其工程实现的核心难题——充气结构形态重构、超大惯量飞行控制、
车载快速部署——已在本文中得到系统性的探索与初步验证。
在新能源汽车加速渗透、个人出行场景不断被重新定义的时代背景下，
"随车常备、自动跟随、双手解放"的飞行雨伞具备成为下一代个人出行配件的技术基础。
本文的工作为这一方向的后续研究与产品化探索奠定了工程基础。